
Machine learning researchers waste hours to days implementing hypotheses only to discover critical feasibility constraints post-implementation. The layer-wise logit extraction example (68.65\% overhead discovered after full implementation) motivated this work: could micro-pilot empirical measurement enable early viability prediction before resource commitment?

Pilot-Driven Viability Gates treats viability assessment as incremental empirical validation across sample scales (10 samples, 100 samples, full dataset) with Bayesian posterior refinement. The framework addresses a methodological gap: ablation studies test hypothesis variations, but no formalized protocol exists for viability gates that identify non-viable hypotheses at micro-pilot stage.

In validation on 32 synthetic ML hypotheses, the framework achieved:

\begin{enumerate}
\item Perfect linear scaling (r=1.000, p<0.0001) between 10-sample micro-pilot and full-scale overhead, validating the extrapolation mechanism in idealized conditions.
\end{enumerate}

\begin{enumerate}
\item 40.91\% Bayesian error reduction (p=0.0003) from Gate 1 to Gate 2, demonstrating that mid-scale observations (100 samples) refine micro-pilot predictions.
\end{enumerate}

\begin{enumerate}
\item 93.3\% Gate 1 classification accuracy (28/30 correct, p=4.$34\times 10^{-7})$, exceeding the 80\% target by 13.3 percentage points and significantly outperforming random baseline (50\%).
\end{enumerate}

\begin{enumerate}
\item 96.2\% recall on non-viable hypotheses (25/26 identified), validating the framework's design priority to minimize false negatives at the cost of occasional false positives.
\end{enumerate}

These results establish proof-of-concept for framework mechanics. However, perfect linearity (r=1.000, k=1.000, CV=0.00\%) reflects synthetic data artifact. Real-world validation must test whether r $\geq$ 0.7 correlation threshold holds under measurement noise, hardware variance, and non-linear scaling effects.

\subsubsection{Future Directions}

Validation results motivate prioritized research directions grounded in observed limitations:

\textbf{High Priority:}

\textbf{Real Corpus Validation:} Motivated by perfect linearity artifact (r=1.000). Collect 30+ published papers from Papers with Code, NeurIPS/ICML/ICLR (2020-2024) reporting both micro-pilot and full-scale overhead. Test whether r $\geq$ 0.7 threshold holds with real measurement noise and hardware variance. Expected outcome: r=0.7-0.9 (lower than synthetic but still above threshold). This establishes external validity beyond synthetic proof-of-concept.

\textbf{Prospective User Study:} Motivated by untested user compliance. Recruit 30 ML researchers, assign to framework vs control groups, measure stop rate at Gate 1 (target $\geq 60$\%), time savings (target $\geq 3$ hours per non-viable hypothesis), and false negative rate (<10\%). This validates framework adoption as decision-making tool, not just predictive accuracy.

\textbf{Multi-Threshold Validation:} Motivated by single threshold limitation (10\% only). Test classification at 10\%, 30\%, 50\%, 100\%, 200\% thresholds. Measure accuracy and confusion matrix at each threshold. Identify threshold range where accuracy $\geq 80$\% (framework applicability scope). Expected pattern: Accuracy decreases as threshold increases (viable/non-viable distribution shifts from 87\% non-viable at 10\% to 50\% at 50\% threshold).

\textbf{Medium Priority:}

\textbf{Non-Linear Scaling Robustness:} Motivated by perfect linear scaling assumption. Collect 20 hypotheses with known non-linear patterns (memory bottlenecks, I/O overhead). Test whether r < 0.7 (correlation fails). If correlation breaks, develop memory profiling extension for Gate 1. Expected outcome: Non-linear corpus r=0.4-0.6 (below threshold), triggering framework extension.

\textbf{Per-Type Scaling Calibration:} Motivated by untested Assumption A3 (k generalizes). After real corpus validation, compute k\_attention, k\_gradient, k\_regularization, k\_normalization. Test whether CV < 30\% (single global k sufficient). If CV > 30\%, develop per-type lookup table. Expected outcome: Real corpus CV=10-20\% (moderate variance but below threshold).

Each future direction addresses a specific limitation, prioritized by impact on framework validity (real corpus) and adoption (user study) over mechanism refinement.

\subsubsection{Closing Perspective}

As machine learning research continues to scale in computational demands and deployment constraints, the gap between hypothesis generation and resource-constrained validation widens. This framework provides a systematic early-stop protocol, filling a methodological need identified by the prevalence of post-implementation feasibility discoveries. While synthetic validation establishes that framework mechanics work as designed, the path from proof-of-concept to practical adoption requires real-world validation on published corpora and prospective user studies.

This work encourages the ML community to formalize viability assessment as a distinct methodological stage---neither hypothesis generation nor performance optimization, but a principled empirical checkpoint that asks ''should we implement this at all?'' before committing research effort. The 93.3\% accuracy achieved in synthetic validation suggests this question can be answered reliably at micro-pilot scale, saving resources for hypotheses that genuinely warrant full-scale investigation.
